\documentclass[aspectratio=169,10pt]{beamer}
\usetheme{Madrid}
\usecolortheme{default}
\usepackage{graphicx}
\usepackage{array}
\setbeamertemplate{navigation symbols}{}
\newcommand{\todo}[1]{\textbf{[TODO: #1]}}

\title[DFF-Adapter Reproduction]{Reproducing DFF-Adapter: Fine-Grained DINOv2 Tuning with Dual Supervision for Face Forgery Detection}
\author[Sudhiksha Mattewada, Moguluri Sai Siddhardha]{Sudhiksha Mattewada-B24CS027, Moguluri Sai Siddhardha-B24DS018 }
\date{}

\begin{document}

\begin{frame}\titlepage\end{frame}

\begin{frame}{Problem and motivation}
\begin{itemize}
\item Face-forgery detectors score $\approx$99\% on their training benchmark but often fail on unseen manipulations
\item Question: does a \textbf{parameter-efficient} DINOv2 detector reproduce both strong FF++ accuracy \emph{and} cross-dataset generalization?
\item Paper: Fine-Grained DINO Tuning with Dual Supervision for Face Forgery Detection (2511.12107v1)
\item Criterion: reproduce paper results (own re-implementation); extension [attempted]: label-smoothed type loss
\end{itemize}
\end{frame}

\begin{frame}{Baseline paper: DFF-Adapter}
\begin{itemize}
\item Frozen DINOv2 (paper: registers variant, \textbf{Large}), 3.5M trainable parameters
\item Multi-head LoRA-expert adapters on query, value and output projections of \textbf{every} block (prior work: last block or last layers only)
\item Rank 16, $\alpha=32$, 6 experts, 4 heads: 1 shared head + 3 task heads (top-3 routing); FAMHR router + SETF fusion
\item Dual supervision during training; inference uses fused authenticity + shared branches only
\item Reported: best overall on five benchmarks, strong cross-dataset / cross-manipulation results
\end{itemize}
\end{frame}

\begin{frame}{Our approach}
\centering
\includegraphics[height=0.66\textheight]{Pipeline.png}
\vspace{2pt}
\begin{itemize}\footnotesize
\item $\mathcal{L}=10\cdot\mathrm{BCE}(\text{auth})+2\cdot\mathrm{CE}(\text{type})$; two passes per batch; Small: 667{,}374 trainable parameters
\item Extension: type targets $(1-\epsilon)\,e_y+\epsilon/K$, $\epsilon=0.1$, type loss only
\end{itemize}
\end{frame}

\begin{frame}{Data and Experimental Setup}
\footnotesize
\textbf{Training Data and Preprocessing}
\begin{itemize}
    \item \textbf{Dataset:} FaceForensics++ (FF++) c23
    \item \textbf{Size:} 5,000 videos across 5 classes (original + 4 manipulations)
    \item \textbf{Preprocessing:} MTCNN face detection and cropping to $224\times224$
    \item \textbf{Frame extraction:} 32 frames per video
\end{itemize}
\vspace{2pt}
\textbf{Cross-Dataset Evaluation}
\begin{itemize}
    \item \textbf{Celeb-DF-v1:} 100 videos, 3,200 frames
    \item \textbf{Celeb-DF-v2:} 518 videos, 16,576 frames
    \item \textbf{DFDC subset:} 392 videos (76 real, 316 fake), 12,544 frames
\end{itemize}
\vspace{2pt}
\textbf{Experimental Configuration}
\begin{itemize}
    \item \textbf{Backbone:} Frozen DINOv2-Small / Base
    \item \textbf{Training:} 12 / 30 epochs; batch size 24
    \item \textbf{Frames per video:} training-10, validation-8
    \item \textbf{Optimizer settings:} Learning rate $10^{-4}$; weight decay 0.01
    \item \textbf{Environment:} Random seed 42; Kaggle NVIDIA T4
    \item \textbf{Metrics:} Video-level AUC (primary), EER and frame-level AUC
\end{itemize}
\vspace{4pt}
\centering
\scriptsize
\textit{Paper configuration: DINOv2-Large, 50 epochs, learning rate $2\times10^{-4}$, weight decay $10^{-5}$, RTX 4090.}
\end{frame}

\begin{frame}{Results: Cross-Dataset Generalization}
\scriptsize
\begin{columns}[T,onlytextwidth]
\column{0.5\textwidth}
\centering
\textbf{Video-level AUC (\%)}
\vspace{3pt}
\begin{tabular}{|l|c|c|c|c|}
\hline
Model (Ep.) & FF++ & CDF-v1 & CDF-v2 & DFDC\\\hline
Small (12) & 99.08 & 81.20 & 85.02 & --\\
\textbf{Small (30)} & \textbf{99.08} & \textbf{82.98} & \textbf{86.25} & \textbf{88.03}\\
Small + LS (30) & 99.05 & 81.20 & 84.74 & 87.45\\
Base (12) & 99.08 & 82.43 & 81.51 & --\\
Base (30) & 99.16 & 81.71 & 84.22 & 84.00\\
Paper Large (50) & 99.56 & 96.14 & 95.26 & 89.96\\\hline
\end{tabular}
\column{0.5\textwidth}
\centering
\textbf{Equal Error Rate (\%)}
\vspace{3pt}
\begin{tabular}{|l|c|c|c|c|}
\hline
Model (Ep.) & FF++ & CDF-v1 & CDF-v2 & DFDC\\\hline
Small (12) & -- & 30.65 & 22.94 & --\\
\textbf{Small (30)} & \textbf{3.57} & \textbf{26.32} & \textbf{22.35} & \textbf{19.74}\\
Small + LS (30) & 2.86 & 26.32 & 25.00 & 19.94\\
Base (12) & -- & 25.81 & 26.18 & --\\
Base (30) & 3.57 & 28.95 & 23.82 & 23.68\\
Paper Large (50) & -- & n/r & 12.65 & 17.96\\\hline
\end{tabular}
\end{columns}
\vspace{8pt}
\hrule
\vspace{6pt}
\footnotesize
\begin{itemize}
\item \textbf{Best Performer:} Small (30 epochs) achieves our best overall cross-dataset performance.
\item \textbf{Label Smoothing:} Yields the lowest FF++ EER (2.86\%) but does not improve cross-dataset AUC.
\item \textbf{Generalization Gap:} Remains weaker on Celeb-DF, while our DFDC subset results are closer to the paper.
\item \textbf{Consistency:} EERs correspond to the model checkpoint with the highest validation video-AUC.
\end{itemize}
\vspace{4pt}
\centering
\scriptsize
\textit{Note: FF++ shows our validation metrics versus the paper's test metrics; the splits differ.}
\end{frame}

\begin{frame}{Experimental Comparisons: Training, Backbone, and Extensions}
\scriptsize
\begin{columns}[T,onlytextwidth]
\column{0.55\textwidth}
\centering
\textbf{12 $\to$ 30 epochs} ($\Delta$ pp)\\[2pt]
\begin{tabular}{|l||c|c|c|c|}
\hline
Model & AUC v1 & EER v1 & AUC v2 & EER v2\\\hline
Small & $+1.78$ & $-4.33$ & $+1.23$ & $-0.59$\\
Base & $-0.72$ & $+3.14$ & $+2.71$ & $-2.36$\\\hline
\end{tabular}
\column{0.43\textwidth}
\centering
\textbf{Base $-$ Small, 30 ep.} ($\Delta$ AUC pp)\\[2pt]
\begin{tabular}{|l||c|}
\hline
Celeb-DF-v1 & $-1.27$\\
Celeb-DF-v2 & $-2.03$\\
DFDC & $-4.03$\\
FF++ val & $+0.08$\\\hline
\end{tabular}
\end{columns}
\vspace{6pt}
\centering
\textbf{Label Smoothing as a Limitation-Focused Extension (LS $-$ baseline)}\\[2pt]
\begin{tabular}{|l||c|c|c|}
\hline
Small, 30 ep. & $\Delta$ video AUC & $\Delta$ EER & $\Delta$ frame AUC\\\hline
Celeb-DF-v1 & $-1.78$ & $0.00$ & n/a\\
Celeb-DF-v2 & $-1.51$ & $+2.65$ & n/a\\
DFDC subset & $-0.58$ & $+0.20$ & $+0.91$\\\hline
\end{tabular}
\vspace{8pt}
\hrule
\vspace{6pt}
\footnotesize
\begin{itemize}
\item \textbf{Training Duration:} Longer training improves Small's cross-dataset results, while Base shows mixed effects.
\item \textbf{Backbone Size:} Base gives only a marginally higher FF++ validation AUC (+0.08 pp), but performs worse on the unseen datasets.
\item \textbf{Label Smoothing:} No video AUC gain on any unseen set; DFDC frame AUC improved slightly (+0.91 pp).
\item \textbf{Architectural Focus:} The paper reports benefits from its router and shared task fusion; our current experiments focus on backbone size, training duration, and label smoothing.
\end{itemize}
\end{frame}

\begin{frame}{Results at a glance}
\centering
\includegraphics[height=0.72\textheight]{fig_auc}
\end{frame}

\begin{frame}{Discussion}
\begin{columns}[T]
\column{0.5\textwidth}
\includegraphics[width=\textwidth]{fig_gap}
\column{0.5\textwidth}
\begin{itemize}
\item In-domain accuracy $\ne$ robustness: $-16.10$ pp (v1), $-12.8$ pp (v2), $-11.05$ pp (DFDC) from FF++
\item Gap to paper is concentrated on Celeb-DF; causes (backbone, recipe, protocol, implementation) are not separable here
\item Bigger backbone and longer training do not guarantee better generalization
\end{itemize}
\end{columns}
\end{frame}

\begin{frame}{Conclusions and Path Forward}
\footnotesize
\textbf{Conclusions}
\begin{itemize}
\item \textbf{Generalization:} FF++ AUC matches closely, and the DFDC subset is within 2 pp (88.03\% vs.\ 89.96\%). However, the Celeb-DF generalization gap was not reproduced.
\item \textbf{Configurations:} Small (30 epochs) performed best. Using a larger backbone (Base) or adding label smoothing did not improve cross-dataset robustness.
\item \textbf{Takeaway:} Source-dataset performance alone is insufficient to judge a deepfake detector.
\end{itemize}
\vspace{6pt}
\textbf{Limitations $\to$ Next Steps}
\begin{itemize}
\item \textbf{Evaluation Protocol:} We used unverified splits and mixed checkpoints, making direct comparisons difficult. \textit{Next step:} Evaluate a single, consistent checkpoint across all datasets using strictly verified splits.
\item \textbf{Setup \& Statistical Noise:} We used smaller backbones (Small/Base) and a single seed, meaning 1--2 pp differences are within noise. \textit{Next step:} Test DINOv2-Large with the paper's exact recipe across $\ge$3 seeds.
\item \textbf{Ablation Clarity:} Current experiments show \emph{what} changes, not \emph{why}. \textit{Next step:} Run strict one-factor ablations to isolate the adapter, type loss, and label smoothing.
\end{itemize}
\end{frame}

\begin{frame}{References}
\footnotesize
\begin{thebibliography}{9}
\bibitem{a} T. Zhang et al., ``Fine-Grained DINO Tuning with Dual Supervision for Face Forgery Detection,'' AAAI, vol.~40, no.~15, pp.~12780--12788, 2026.
\bibitem{b} M. Oquab et al., ``DINOv2: Learning Robust Visual Features without Supervision,'' TMLR, 2024.
\bibitem{c} E. J. Hu et al., ``LoRA: Low-Rank Adaptation of Large Language Models,'' ICLR, 2022.
\bibitem{d} A. R\"ossler et al., ``FaceForensics++,'' ICCV, 2019.
\bibitem{e} Y. Li et al., ``Celeb-DF,'' CVPR, 2020.
\bibitem{f} B. Dolhansky et al., ``The DeepFake Detection Challenge (DFDC) Dataset,'' arXiv:2006.07397, 2020.
\end{thebibliography}
\end{frame}

\end{document}
